\documentclass[11pt]{article}
\usepackage[T1]{fontenc}
\usepackage{lmodern,microtype,amsmath,amssymb,amsthm,mathtools}
\usepackage[margin=1in]{geometry}
\usepackage[colorlinks=true,linkcolor=blue,citecolor=blue,urlcolor=blue]{hyperref}
\newtheorem{theorem}{Theorem}[section]
\newtheorem{lemma}[theorem]{Lemma}
\newtheorem{proposition}[theorem]{Proposition}
\newcommand{\R}{\mathbb R}\newcommand{\C}{\mathbb C}
\newcommand{\dom}{\operatorname{dom}}\newcommand{\ran}{\operatorname{ran}}
\newcommand{\Sc}{\mathcal S}\newcommand{\ip}[2]{\langle #1,#2\rangle}
\title{Closed differentiation and signed adjoints\\for Schwartz densities with zeros}
\author{Project Stardust\\\small A classical formalization companion}
\date{October 29, 2026}
\begin{document}
\maketitle
\begin{abstract}
For every real Schwartz density $K$, differentiation on Schwartz
functions is densely closable in $L^2(|K|\,dx)$, without a
square-integrable logarithmic derivative. We give a direct
distributional construction using bounded coefficients and a
weight-squared embedding, keeping the original graph closure
distinct from its closed weak extension. Explicit total adjoint
probes and multiplication by the sign of $K$ identify the exact
domain of the signed adjoint. Orthogonal graph projection constructs
the positive Hilbert energy resolvent, its exact minimization gap
and the full nonnegative self-adjoint energy operator. The
logarithmic derivative is always absolutely integrable against
$|K|\,dx$ but fails to be square-integrable at a simple zero.
A normalized sign-changing polynomial-Gaussian example exhibits
this failure while retaining all closed operators. An integrating
factor also excludes polynomial first-order Stein weights for
sign-changing kernels. The mechanisms are classical weighted
operator and indefinite-inner-product theory; no new general
closability or spectral principle is claimed. A separate Gaussian/Dirac
mixture shows that full support does not replace the density hypothesis.
\end{abstract}

\section{The original spaces and conclusions}
Fix a real Schwartz function $K\in\Sc(\R;\R)$ and write
$w=|K|$, $\mu(dx)=w(x)\,dx$. No positivity, nonvanishing,
normalization or finite-zero hypothesis on $K$ is imposed.
The measure $\mu$ is finite. Throughout,
\[
 H=L^2(\mu;\C),\qquad
 \ip uv=\operatorname{Re}\int\overline{u}v\,d\mu
\]
is regarded as a \emph{real} Hilbert space. All adjoints and
self-adjointness assertions below refer to this real structure,
matching the formal sources. Distribution tests are complex Schwartz
functions with complex bilinear distribution evaluation.
The zero-density case gives the zero Hilbert space and is included.

Let $D_0$ take the $H$ class of a Schwartz function $f$ to the
class of $f'$. Its representative independence will be proved,
rather than presumed. The principal result is the following.
\begin{theorem}\label{thm:main}
The operator $D_0$ is densely defined and closable. Its graph
closure $D$ is densely defined and closed. Set $J u=\sigma u$,
where $\sigma=1$ on $\{K\geq0\}$ and $\sigma=-1$ on $\{K<0\}$.
The signed adjoint is the densely defined closed operator
\[
 D^+=JD^*J,\qquad
 \dom D^+=\{v\in H:Jv\in\dom D^*\}.
\]
For the nondegenerate signed real form
$[u,v]=\operatorname{Re}\int\overline u vK\,dx$, it satisfies
\[
 [Du,v]=[u,D^+v]\quad(u\in\dom D, v\in\dom D^+),
\]
and this pairing condition characterizes its whole domain and value.
The Hilbert energy operator
\[
 A=D^*D,\qquad \dom A=\{u\in\dom D:Du\in\dom D^*\}
\]
is densely defined, closed, nonnegative and self-adjoint. Its
resolvent $R=(I+A)^{-1}$ is a bounded symmetric positive contraction,
injective with dense range. For every $x\in H$ and $u\in\dom D$,
\begin{align}\label{eq:gap}
 \|x-u\|^2+\|Du\|^2
 ={}&\|x-Rx\|^2+\|DRx\|^2\notag\\
 &+\|u-Rx\|^2+\|Du-DRx\|^2.
\end{align}
Thus $Rx$ is the unique minimizer of the displayed energy.
\end{theorem}
The positive operator $A$ uses the \emph{Hilbert} adjoint $D^*$,
not the signed adjoint $D^+$. Positivity means its Hilbert quadratic
form is nonnegative; it is not an assertion of order preservation
or a probabilistic transition kernel. Neither operator is identified
with a time generator of a signed pseudoprocess.

\section{A bounded-coefficient distributional rigging}
For $u\in H$ and $\varphi\in\Sc(\R;\C)$ define
\[
 Iu(\varphi)=\int\varphi u\,d\mu,\qquad
 Eu=I(wu),\qquad Fu=I(bu),
\]
where
\begin{equation}\label{eq:boundeddrift}
 b(x)=2\sigma(x)K'(x),\qquad w(x)b(x)=2K(x)K'(x).
\end{equation}
Both $w$ and $b$ are bounded measurable functions, since $K,K'$
are Schwartz. In particular these formulas define continuous
real-linear maps from $H$ into the complex tempered distributions.
For example Cauchy--Schwarz gives
\[
 |Eu(\varphi)|\leq\|u\|_H\|w\varphi\|_H,
 \qquad |Fu(\varphi)|\leq\|u\|_H\|b\varphi\|_H,
\]
and each second factor is at most a fixed constant times
$\sup|\varphi|$. Thus every evaluation is continuous in $u$;
the bounds also give the usual distribution topology continuity.
No quotient by $K$ appears.

It is important to retain the extra factor in $E$:
\begin{equation}\label{eq:weightsquared}
 Eu(\varphi)=\int\varphi u\,w^2\,dx
             =\int\varphi u\,K^2\,dx.
\end{equation}
This is not the distribution represented by $u$ against Lebesgue
measure, which need not exist locally across zeros. Its density
is integrable, because
$\int|u|w^2\leq\|u\|_H(\int w^3)^{1/2}<\infty$.
The map $E$ is injective. Indeed $Eu=0$ implies $u w^2=0$
Lebesgue almost everywhere by separation of locally integrable
densities by smooth compactly supported tests. One elementary
justification of that separation is to convolve locally with a
compactly supported smooth approximate identity: all convolutions
are zero on smaller intervals, and converge to the density in
local $L^1$. Hence the density is zero there, and a countable
interval covering gives the assertion on $\R$. Since $w>0$
$\mu$-almost everywhere, it follows that $u=0$ in $H$.

For a Schwartz input $f$, integration by parts gives the exact
identity
\begin{equation}\label{eq:coreweak}
 E(f')=(Ef)'-Ff.
\end{equation}
To check all signs, evaluate the right side on $\varphi$. It is
$-\int\varphi'fK^2-\int2\varphi fKK'$, which equals
$\int\varphi f'K^2$. Every term is absolutely integrable, and
the boundary terms vanish because the factors are Schwartz.
This uses the smooth weight $K^2$, even though $|K|$ need not
be differentiable at a transverse zero.

\section{The weak graph and the original derivative closure}
Define
\begin{equation}\label{eq:weakgraph}
 \mathcal G=\{(u,v)\in H\times H:Ev=(Eu)'-Fu\}.
\end{equation}
It is a closed real linear subspace: for a convergent sequence
$(u_n,v_n)$ in it, continuity of each test evaluation passes the
identity to the limit. If $(0,v)\in\mathcal G$, injectivity of
$E$ gives $v=0$. Thus $\mathcal G$ is the graph of a closed
partially defined operator $D_{\mathrm w}$.

Equation \eqref{eq:coreweak} shows that every pair $(f,f')$ of
Schwartz classes lies in this graph. In particular if a Schwartz
$f$ is zero as an $H$ class, then its derivative class is zero:
$D_0$ is well-defined on the range of Schwartz functions in $H$.
This range is dense. For clarity, finite Radon measure regularity
first approximates indicators of measurable sets in $L^2$ by
continuous compactly supported functions: choose compact and open
inner/outer approximations with small measure difference and a
continuous cutoff between them. Truncation and simple-function
approximation reduce a general $L^2$ function to these indicators.
Finally mollification approximates a continuous compactly supported
function uniformly by smooth compactly supported functions, hence
in $L^2(\mu)$ since $\mu$ is finite. Such functions are Schwartz.
The same argument applies to real and imaginary parts.

Since $D_0\subset D_{\mathrm w}$ and the latter has closed
single-valued graph, $D_0$ is closable. Explicitly if $f_n\to0$
and $f_n'\to v$ in $H$, their graph pairs tend to $(0,v)$ in
$\mathcal G$, so $v=0$. The closure of the core graph therefore
defines a densely defined closed operator $D$ with
\begin{equation}\label{eq:minimalweak}
 D_0\subset D\subset D_{\mathrm w}.
\end{equation}
This is an inclusion, not a proved identification of the minimal
graph closure with every distributional solution of
\eqref{eq:weakgraph}. The Schwartz graph is a core for $D$ by
its definition. No different maximal weighted Sobolev domain or
condition at individual zeros is substituted for it.

The nonnegative core form $\kappa\|D_0f\|^2$ is closable for
every $\kappa\geq0$. For $\kappa>0$, a sequence tending to zero
whose derivative energies are Cauchy has derivatives Cauchy in
$H$, hence convergent by completeness, and the preceding graph
test makes their limit zero. For $\kappa=0$ the assertion is
immediate. This is positive \emph{absolute-density} energy, not
an energy obtained by integrating a square against the signed law.

\section{Explicit adjoint probes and their totality}
The real Hilbert adjoint of $D$ is defined by
\[
 v\in\dom D^*\ \Longleftrightarrow\
 \exists z\in H\ \forall u\in\dom D:\ \ip v{Du}=\ip z u,
 \qquad D^*v=z.
\]
Density of $\dom D$ makes $z$ unique. Its graph is closed by
passing these inner-product identities to limits. For every
complex Schwartz test $\varphi$ define explicit $H$ vectors
\begin{equation}\label{eq:probes}
 p_\varphi=w\overline\varphi,\qquad
 q_\varphi=-w\overline{\varphi'}-b\overline\varphi.
\end{equation}
Boundedness of the coefficients and finiteness of $\mu$ put
both in $H$. Their normalization is
\[
 \ip{p_\varphi}u=\operatorname{Re}Eu(\varphi),\quad
 \ip{q_\varphi}u=-\operatorname{Re}Eu(\varphi')
                       -\operatorname{Re}Fu(\varphi).
\]
For $(u,v)\in\mathcal G$, \eqref{eq:weakgraph} therefore gives
$\ip{p_\varphi}v=\ip{q_\varphi}u$. Applied to $v=Du$, this proves
\begin{equation}\label{eq:adjointprobe}
 p_\varphi\in\dom D^*,\qquad D^*p_\varphi=q_\varphi.
\end{equation}
The formal development defines these vectors by real Riesz
representation; \eqref{eq:probes} is the direct coordinate
realization verified by the two displayed pairings. No division
by a vanishing density is needed.

Their real linear span is dense. If $u$ is orthogonal to every
$p_\varphi$, then $\operatorname{Re}Eu(\varphi)=0$ for every
$\varphi$. Testing $i\varphi$ gives the imaginary part as well,
so $Eu=0$ and injectivity gives $u=0$. The closed span has zero
orthogonal complement and is therefore $H$. In particular
$\dom D^*$ is dense. These probes are total in $H$; they have
not been asserted to be a graph core for $D^*$.

For completeness the bidual is the original $D$. A pair $(u,v)$
is in its closed graph precisely when it is orthogonal to every
vector of its graph's orthogonal complement. The latter consists
of $(-D^*z,z)$ for $z\in\dom D^*$: orthogonality to $(f,Df)$
is exactly the adjoint definition. Orthogonality of $(u,v)$ to
all these vectors is $\ip vz=\ip u{D^*z}$, the definition of
$(u,v)$ in the graph of $D^{**}$. Thus $D^{**}=D$, without
replacing the original domain by $\dom D_{\mathrm w}$.

\section{The signed form and its full adjoint domain}
Multiplication by $\sigma$ gives a bounded real-linear involution
$J$: $J^2=I$, $\|Ju\|=\|u\|$, and
$\ip{Ju}v=\ip u{Jv}$. These follow pointwise from
$\sigma^2=1$ and the fact that $\sigma$ is real. The signed form
\[
 [u,v]=\ip{Ju}v=\operatorname{Re}\int\overline u vK\,dx
\]
is continuous and symmetric, and $[u,Ju]=\|u\|^2$. It is
nondegenerate even if $[u,u]=0$ for nonzero $u$; isotropy is
different from degeneracy. The Hilbert space decomposes into
the closed functions supported on $\{K>0\}$ and $\{K<0\}$,
where this form is respectively positive and negative. The zero
set has zero $\mu$ measure. This is the familiar real indefinite
Hilbert-space construction by a fundamental symmetry.

Suppose $v,z\in H$. The identity $[Du,v]=[u,z]$ for every
$u\in\dom D$ is equivalent, using symmetry of $J$, to
\[
 \ip{Du}{Jv}=\ip u{Jz}\quad\hbox{for every }u\in\dom D.
\]
By the Hilbert adjoint definition this holds if and only if
$Jv\in\dom D^*$ and $D^*Jv=Jz$. Multiplying by $J$ proves
the exact domain and formula $D^+v=JD^*Jv$ in
Theorem~\ref{thm:main}. Its graph is the inverse image of the
closed graph of $D^*$ under the continuous invertible map
$(v,z)\mapsto(Jv,Jz)$, so it is closed. Its domain is
$J(\dom D^*)$, dense because $J$ is a Hilbert isometry onto
$H$ and the adjoint domain is dense.

In particular $Jp_\varphi$ is in the signed-adjoint domain with
$D^+(Jp_\varphi)=Jq_\varphi$; these are actual dense test-generated
vectors. The signed adjoint is not obtained by presuming a formal
logarithmic derivative acts on every Schwartz input. Behavior at
zeros must satisfy its actual transported Hilbert-adjoint domain.

\section{Graph projection and the exact energy minimizer}
The graph of $D$ is a closed subspace of the product Hilbert
space $H\oplus H$ with norm squared $\|u\|^2+\|v\|^2$.
Project $(x,0)$ orthogonally onto it and write the projection
as $(Rx,Gx)$. Thus $Rx\in\dom D$ and $Gx=DRx$.
Orthogonal projection and the two coordinate maps are bounded
linear, so $R,G$ are bounded real-linear maps. Orthogonality
of the residual gives, for every $u\in\dom D$,
\begin{equation}\label{eq:variational}
 \ip{Rx}u+\ip{Gx}{Du}=\ip x u.
\end{equation}
The adjoint definition now yields
\[
 Gx\in\dom D^*,\qquad D^*Gx=x-Rx.
\]
Hence $Rx\in\dom D^*D$ and $(I+D^*D)Rx=x$.
Taking $u=Rx$ in \eqref{eq:variational} gives
\begin{equation}\label{eq:energyidentity}
 \|Rx\|^2+\|Gx\|^2=\ip x{Rx}\geq0.
\end{equation}
Projection is contractive, so $\|Rx\|\leq\|x\|$.
Self-adjointness of the bounded graph projection also gives
$\ip{Rx}y=\ip x{Ry}$ by testing the two product inputs $(x,0)$
and $(y,0)$. Thus $R$ is symmetric and positive in the Hilbert
quadratic-form sense.

If $Rx=0$, single-valuedness of the graph gives $Gx=0$.
Equation \eqref{eq:variational} then says $x$ is orthogonal to
the dense domain of $D$, so $x=0$. Thus $R$ is injective.
If $y$ is orthogonal to its range, symmetry implies
$\ip x{Ry}=0$ for every $x$, hence $Ry=0$ and $y=0$.
The range of $R$ is consequently dense.

The orthogonal residual $(x-Rx,-Gx)$ is perpendicular to
$(u-Rx,Du-Gx)$. Pythagoras applied to their difference gives
the exact gap \eqref{eq:gap}. The last two terms there are
nonnegative, so $Rx$ minimizes $\|x-u\|^2+\|Du\|^2$.
If equality holds for another input, $\|u-Rx\|^2=0$ gives
uniqueness. Both its domain membership and its derivative value
have been established before writing this minimization statement.

\section{The full Hilbert energy operator}
We finish the operator part without invoking an unspecified
maximal differential expression. If $u\in\dom D$ and
$v=Du\in\dom D^*$, set $z=D^*v$ and $x=u+z$.
For every $a\in\dom D$, the adjoint identity gives
$\ip x a=\ip u a+\ip v{Da}$. Comparing with
\eqref{eq:variational} and testing the difference against
$a=u-Rx$ yields
\[
 \|u-Rx\|^2+\|v-Gx\|^2=0.
\]
Thus $u=Rx$ and $z=x-Rx$. Conversely every $Rx$ has exactly
this factorization by the previous section. Therefore
\begin{equation}\label{eq:factorization}
 \dom A=\ran R,\qquad A(Rx)=x-Rx,\qquad A=D^*D
\end{equation}
with the precise product domain in Theorem~\ref{thm:main}.
Injectivity of $R$ makes this single-valued and density of its
range makes $A$ densely defined. Its quadratic form is
$\ip{Au}u=\|Du\|^2\geq0$.

Symmetry follows either from $D^*D$ or directly from symmetry
of $R$: $\ip{x-Rx}{Ry}=\ip{Rx}{y-Ry}$.
To prove full self-adjointness, let $z\in\dom A^*$ with
$A^*z=y$. Testing against $Rx$ for arbitrary $x\in H$ gives
\[
 \ip{x-Rx}z=\ip{Rx}y.
\]
Symmetry of $R$ makes this $\ip x{z-Rz-Ry}=0$ for all $x$.
Hence $z=R(z+y)$, so $z\in\ran R=\dom A$ and
$Az=z+y-R(z+y)=y$. Thus $A^*\subseteq A$; symmetry gives
the reverse inclusion. An adjoint is closed by its defining
pairings, so $A=A^*$ is also closed. Equation
\eqref{eq:factorization} proves that $I+A$ is onto and its
inverse is exactly the constructed $R$. This completes every
operator assertion of Theorem~\ref{thm:main}.

The weak distributional factorization retained by these actual
solutions is
\[
 E(DRx)=(E(Rx))'-F(Rx),\qquad D^*(DRx)=x-Rx.
\]
It does not identify $A$ with an unqualified pointwise
Sturm--Liouville expression or identify an energy-semigroup
law with the signed density $K\,dx$.

\section{An integrable score, simple zeros and a normalized example}
Define $s_K=-K'/K$ where $K\ne0$ and set it to zero where
$K=0$, the division convention of the source. The cancellation
$Ks_K=-K'$ holds Lebesgue almost everywhere, including the
zero set up to a null subset. To see the last qualification,
every zero where $K'\ne0$ is isolated among zeros, by the
definition of a nonzero derivative. A discrete subset of $\R$
is countable: choose for each isolated point a rational interval
containing no other point, and use the countable rational basis.
Thus $K'=0$ at almost every zero of $K$.
Consequently
\begin{equation}\label{eq:scoreL1}
 \int |s_K|\,d\mu=\int|K'|\,dx<\infty.
\end{equation}
For every Schwartz test $f$, ordinary integration by parts gives
the actual signed Stein identity
\[
 \int f'K\,dx=\int f s_K K\,dx.
\]
The right side is absolutely integrable since $f$ is bounded and
\eqref{eq:scoreL1} holds. Existence of this signed integral does
not imply the score belongs to $H$.

Suppose $K(0)=0$ and $K'(0)\ne0$. Continuity of $K'$ and
the derivative definition give $c,C,\delta>0$ such that, for
$0<x<\delta$,
\[
 |K'(x)|\geq c,\qquad 0<|K(x)|\leq Cx.
\]
Indeed the first inequality follows by staying within half of
$|K'(0)|$; the quotient $K(x)/x\to K'(0)$ gives both
nonvanishing and the upper bound for sufficiently small $x$.
Then
\[
 \int|s_K|^2\,d\mu
 \geq\int_0^\delta\frac{|K'(x)|^2}{|K(x)|}\,dx
 \geq\frac{c^2}{C}\int_0^\delta\frac{dx}{x}=\infty.
\]
Thus a transverse zero obstructs $L^2$ score integrability,
while all of Theorem~\ref{thm:main} still holds.

For a concrete normalized signed density set
\[
 K_0(x)=\frac{(x+x^2)e^{-x^2}}{M},\qquad
 M=\int_\R(x+x^2)e^{-x^2}\,dx>0.
\]
The odd term integrates to zero by $x\mapsto-x$. The remaining
$x^2e^{-x^2}$ is nonnegative and strictly positive on $(1,2)$,
so its integral is positive and finite. Hence $K_0$ is Schwartz,
has integral one and is strictly negative at $x=-1/2$.
Also $K_0(0)=0$ and $K'_0(0)=1/M\ne0$. It follows that
$s_{K_0}\in L^1(|K_0|\,dx)\setminus L^2(|K_0|\,dx)$.
Its original derivative closure, dense signed adjoint and positive
Hilbert energy resolvent/operator nevertheless exist. This
counterexample does not say the score of every sign-changing
higher-heat kernel has a simple zero; that extra zero hypothesis
has not been proved or assumed here.

\section{Polynomial Stein weights and higher even heat marginals}
There is a different first-order obstruction that needs no
simple-zero hypothesis. Suppose a real polynomial $P$ satisfies
\[
 \int f'K\,dx=\int f P K\,dx
 \quad\hbox{for every real Schwartz }f.
\]
Integration by parts says $\int f(K'+PK)=0$ for all such $f$.
The residual $K'+PK$ is itself Schwartz; choosing $f$ to be
that residual gives its squared integral zero. Continuity then
gives $K'+PK=0$ at every point. If
$Q(x)=\int_0^xP(y)\,dy$, differentiation yields
\[
 (e^{Q(x)}K(x))'=e^{Q(x)}(K'(x)+P(x)K(x))=0.
\]
The expression is constant on $\R$ by the mean value theorem.
Since the exponential is strictly positive, a nonzero such $K$
has one fixed sign. Thus a sign-changing real Schwartz density
cannot have any polynomial first-order Stein weight of this form.
This does not preclude the measurable $L^1$ score above, higher
derivative identities or a calculus on a specified test dual.

For example let $p=2m>2$ and $t>0$, and let $k_t$ be the real
even Schwartz kernel with angular Fourier transform
$\widehat{k_t}(\xi)=e^{-t\xi^p}$. Its mass is one and its
second moment is zero, by differentiating that transform twice
at zero. If $k_t\geq0$ almost everywhere, the nonnegative
function $x^2k_t(x)$ would have zero integral, forcing $k_t=0$
away from zero almost everywhere, contradicting its unit
Lebesgue mass. Thus it has a negative value; positive mass also
gives a positive value. It is sign-changing, so the polynomial
first-order weight is impossible. On the other hand its
absolute-density space has every closed operator constructed
above. The formal library uses Fourier phases $e^{2\pi i\xi x}$;
rescaling angular frequency changes neither zero second moment
nor sign and operator conclusions.

This is a statement about each fixed signed marginal. It gives
neither a finite-variation signed path law nor predictable
representation, and its positive energy is not a signed-square
isometry. The separate completed multichannel test-dual calculus
has different domains and must retain them.


\section{Full support does not replace the density hypothesis}
The preceding theorem uses an absolutely continuous reference
$|K|\,dx$. A different, positive probability measure can have full
support and a well-defined smooth derivative but no closable
extension of that derivative. Let $\gamma$ be the standard real
Gaussian law and put
\[
 \widetilde\mu=\tfrac12\gamma+\tfrac12\delta_0.
\]
It is a finite Radon probability with full support, since every
nonempty open interval has positive Gaussian mass. Its atom at
zero has mass $1/2$. Consider the ordinary real derivative on
$C_c^\infty(\R)$ in the real space $L^2(\widetilde\mu)$.
It is well-defined: if two continuous representatives agree
$\widetilde\mu$-almost everywhere, they agree everywhere. Otherwise
continuity gives a nonempty open interval where their difference
stays nonzero, contradicting full support. Their ordinary
derivatives therefore agree. The domain is dense by finite Radon
regularity and mollification, just as in Section 3.

Choose a smooth compact cutoff $\chi$ equal to one for $|x|\leq1$
and zero for $|x|\geq2$. Such a cutoff is explicit: take
$\eta(s)=e^{-1/s}$ for $s>0$, zero for $s\leq0$, and set
\[
 \chi(x)=\frac{\eta(4-x^2)}{\eta(4-x^2)+\eta(x^2-1)}.
\]
The denominator is everywhere positive. Smoothness at the seams
follows because every derivative of $e^{-1/s}$ is a finite sum
of polynomial powers of $1/s$ times $e^{-1/s}$ and tends to
zero as $s\downarrow0$. Thus $\chi\in C_c^\infty$ with the
specified values. Write $\psi(x)=x\chi(x)$ and
\[
 f_n(x)=\frac{\psi((n+1)x)}{n+1},\qquad
 f_n'(x)=\psi'((n+1)x).
\]
Then $f_n(0)=0$, $f'_n(0)=1$ and
$\sup|f_n|\leq\sup|\psi|/(n+1)\to0$. For each $x\ne0$,
compact support makes $f'_n(x)=0$ eventually, whereas at zero
it remains one. The derivatives are uniformly bounded by
$\sup|\psi'|$. Dominated convergence for the finite measure
therefore gives
\begin{equation}\label{eq:atomicjetlimit}
 f_n\longrightarrow0,\qquad
 f_n'\longrightarrow 1_{\{0\}}
 \quad\hbox{in }L^2(\widetilde\mu),\qquad
 \|1_{\{0\}}\|_2^2=\tfrac12.
\end{equation}
The graph has a nonzero vertical limit. Hence no closable linear
operator agreeing with all these smooth derivative data can exist:
its graph closure would contain both $(0,0)$ and
$(0,1_{\{0\}})$ and would not be single-valued.

The positive derivative form also fails closability. For $\kappa>0$,
\eqref{eq:atomicjetlimit} implies
\[
 \kappa\|f_n'\|_2^2\longrightarrow\kappa/2>0,
 \qquad \kappa\|f_n'-f_m'\|_2^2\longrightarrow0
 \quad(n,m\to\infty),
\]
although $f_n\to0$ in the underlying space. This violates the
defining sequential form-closability test. At $\kappa=0$ the
zero form is trivially closable. Full support and uniqueness of
continuous representatives do not repair this graph obstruction.

For the pure Dirac law the defect occurs earlier: $\psi$ and
the zero function have the same $L^2(\delta_0)$ class since
$\psi(0)=0$, but their derivatives have different classes since
$\psi'(0)=1$. Thus ordinary smooth differentiation does not even
descend to those equivalence classes. The Gaussian mixture restores
representative independence while preserving nonclosability.
These are separate positive-measure examples, not failures of
Theorem~\ref{thm:main} for a Schwartz absolute density.

\section{Classical antecedents and exact source boundary}
Weighted derivative closability is classical. De Cicco--Serra
Cassano~\cite{DCSC}, Section 2.2/condition (12), recall Hamza's
one-dimensional criterion. Continuous $w$ is locally bounded below
where positive, so $1/w$ is locally integrable there. Their
bounded-interval relaxation results do not identify our global
minimal and weak graphs. Derkach--Trunk~\cite{DT}, Section 2.1,
give the classical $T^+=JT^*J$ mechanism. Their coupling and
spectral theorems are not used. Graph projection and $D^*D$ are
classical Hilbert theory; no priority is claimed.

The exact theorem map records frozen declarations at
\begin{center}\texttt{\small bab688bb6140cd6316c80680095a914ea0d346db}.\end{center}
The map also records the actual smooth shrinking jets and mixed
Gaussian/Dirac nonclosability; the explicit cutoff above realizes
the same unit derivative and radius-one/two support properties.
Riesz-probe coordinates and regrouped human proofs are labelled
reconstructions. The original full-build/axiom evidence is
hash-checked; no research source or original domain is changed.
This formalization companion awaits independent referee review.

\begin{thebibliography}{9}
\bibitem{DCSC} V. De Cicco and F. Serra Cassano,
\emph{Relaxation and optimal finiteness domain for degenerate
quadratic functionals---one dimensional case}, primary author
preprint, author-hosted 2024 version.
\url{https://cvgmt.sns.it/media/doc/paper/6038/DCSC-Relaxation_22_1_2024.pdf}.
Section 2.2/condition (12) comparison only.
\bibitem{DT} V. Derkach and C. Trunk,
\emph{Coupling of definitizable operators in Kre\u\i n spaces},
TU Ilmenau preprint M 17/03 (2017).
\url{https://d-nb.info/1125462493/34}.
Section 2.1 comparison only.
\end{thebibliography}
\end{document}
